\uuid{Ks9H}
\titre{Continuité, dérivées partielles et non-différentiabilité en l'origine}
\niveau{L2}
\module{Analyse}
\chapitre{Fonctions de plusieurs variables}
\sousChapitre{Continuité et différentiabilité}
\theme{Continuité, dérivées partielles, non-différentiabilité, coordonnées polaires}
\auteur{Maxime Nguyen}
\datecreate{2026-03-24}
\organisation{amscc}
\difficulte{4}

\contenu{
\texte{
  On donne $f : \mathbb{R}^2 \to \mathbb{R}$ définie par
  \[
    f(x,y) = \begin{cases} \dfrac{x^3 + 2x^2y + 3y^3}{x^2+y^2} & \text{si } (x,y) \neq (0,0), \\ 0 & \text{sinon.} \end{cases}
  \]
  Montrer que $f$ est continue en $(0,0)$, que les dérivées partielles $\dfrac{\partial f}{\partial x}(0,0)$ et $\dfrac{\partial f}{\partial y}(0,0)$ existent, et que $f$ n'est pas différentiable en $(0,0)$.
}
\begin{enumerate}
  \item \question{Continuité en $(0,0)$.}
    \reponse{
      En coordonnées polaires :
      \[
        f(r\cos\theta, r\sin\theta) = r\bigl(\cos^3\theta + 2\cos^2\theta\sin\theta + 3\sin^3\theta\bigr).
      \]
      Donc $|f(r\cos\theta, r\sin\theta)| \leq 6r \to 0$ quand $r \to 0$. $f$ est continue en $(0,0)$.
    }

  \item \question{Existence des dérivées partielles en $(0,0)$.}
    \reponse{
      $f(h,0) = h$ pour $h \neq 0$, donc $\dfrac{\partial f}{\partial x}(0,0) = \displaystyle\lim_{h\to 0}\dfrac{h}{h} = 1$.

      $f(0,h) = 3h$ pour $h \neq 0$, donc $\dfrac{\partial f}{\partial y}(0,0) = \displaystyle\lim_{h\to 0}\dfrac{3h}{h} = 3$.
    }

  \item \question{Non-différentiabilité en $(0,0)$.}
    \reponse{
      Si $f$ était différentiable en $(0,0)$, on aurait
      \[
        \lim_{(x,y)\to(0,0)} \frac{f(x,y) - x - 3y}{\sqrt{x^2+y^2}} = 0.
      \]
      En coordonnées polaires :
      \[
        \frac{f(r\cos\theta, r\sin\theta) - r\cos\theta - 3r\sin\theta}{r} = 2\cos^2\theta\sin\theta,
      \]
      qui dépend de $\theta$. Donc $f$ n'est pas différentiable en $(0,0)$.
    }
\end{enumerate}
}
